\documentclass[12pt,spanish, singlespacing,]{MastersDoctoralThesis}
\usepackage[utf8]{inputenc} 
\usepackage[T1]{fontenc} 
\usepackage[none]{hyphenat}
\usepackage{verbatim}
\usepackage{palatino} 
\usepackage{acronym}
\usepackage{algorithm}
\usepackage{algorithmic}
\usepackage{datetime}
\usepackage{graphicx}
\usepackage{subfig}
\spacing{1.4}
\usepackage[acronym,nomain]{glossaries}
%\usepackage[backend=bibtex,style=authoryear,natbib=true]{biblatex}
%\addbibresource{main.bib}
\usepackage[autostyle=true]{csquotes}

\usepackage{afterpage}
\newcommand\blankpage{%
    \null
    \thispagestyle{empty}%
    \addtocounter{page}{0}%
    \newpage}
    
\usepackage{lipsum}  

\usepackage{tocloft}    
\newcommand{\listequationsname}{Índice de ecuaciones}
\newlistof{myequations}{equ}{\listequationsname}
\newcommand{\myequations}[1]{%
\addcontentsline{equ}{myequations}{\hspace{1.5em} \protect\numberline{\theequation}#1}\par}
\setlength{\cftmyequationsnumwidth}{2.5em}% Width of equation number in List of Equations

    
\usepackage{listings}
\renewcommand{\lstlistingname}{Código}
\renewcommand{\lstlistlistingname}{Índice de algoritmos}



\usepackage{float}
\usepackage{url}
\usepackage[font=footnotesize,labelfont=bf]{caption}
\makeatletter
\g@addto@macro{\UrlBreaks}{\UrlOrds}
\makeatother

\thesistitle{}
\supervisor{Proin Aliquam Vestibulum} 
\examiner{} 
\degree{} 
\author{Mauris Finibus Ante Ut} 
\subject{}
\keywords{} 
\university{\href{http://www.uni.edu.pe/}{Universidad Nacional de Ingeniería}} 
\department{\href{http://fc.uni.edu.pe/fc/index.php/escuelas/ciencia-de-la-computacion}{}} 
\faculty{\href{http://fc.uni.edu.pe/fc/}{}}
\hypersetup{pdftitle=\ttitle} 
\hypersetup{pdfauthor=\authorname}
\hypersetup{pdfkeywords=\keywordnames} 
\sloppy
\decimalpoint
\begin{document}
\frontmatter 
\pagestyle{plain} 
\begin{titlepage}
\begin{center}
%\textsc{\huge \univname}\\[0.5cm]

\begin{figure}[h]
\centering
\includegraphics[width=0.3\textwidth]{Figures/logo_uni.png}
\end{figure}

\textsc{\huge \univname}\\[0.3cm]
\textsc{\Large Facultad de Ciencias}\\[0.2cm]
\textsc{\large Escuela Profesional de Ciencia de la Computación}\\[1.5cm]
% \textsc{\LARGE \textit{T\'itulo del trabajo}}\\[2cm] 
\textsc{\large Proyecto de Tesis 1}\\[1.5cm] 
{\Large \textbf{Titulo del Trabajo de Proyecto de Tesis N}}
{\huge \bfseries \ttitle}\\[1cm] 
% Thesis title
%\bigskip
%\HRule \\[0.8cm] % Horizontal line

%\begin{minipage}{1.5\textwidth}
%\begin{flushleft} \large
\bigskip
\bigskip
\large\emph{Autor:\\}
{\authorname}\\ 
\large\emph{Asesor:\\} 
{\supname} 
%\end{flushleft}
%\end{minipage}
%\\[2cm]
\\[1cm]
{\large \today}\\[4cm] 
 
\vfill
\end{center}
\end{titlepage}
\afterpage{\blankpage}
%\cleardoublepage

%\renewcommand{\abstractname}{Abstract}
\begin{abstract}
\addchaptertocentry{\abstractname}
%%%%%%%%%%%%%%%%%%%%%%%%% AQUI %%%%%%%%%%%%%%%%%%%%%%%%%%%%%%5
%Poner un resumen del trabajo de no más de 1 página

%El presente trabajo de investigación se centra en las técnicas empleadas para el reconocimiento de voz. 

In cursus dignissim purus a efficitur. Nullam et sem nec tellus rutrum sagittis ut nec augue. Integer suscipit aliquam metus eu fermentum. Vivamus id dolor risus. Suspendisse potenti. Morbi finibus varius lectus eu mattis. Vivamus sapien lacus, finibus nec tortor vel, ultrices efficitur erat. Praesent finibus vitae nisl id iaculis. Praesent congue mollis libero, sit amet luctus ipsum porta non. Praesent pretium et sapien eget hendrerit. Vestibulum feugiat diam sit amet felis feugiat, sed ultricies nulla consectetur. Nam condimentum malesuada orci nec dictum. Suspendisse elementum tincidunt quam, at facilisis metus tempor id. Sed velit tortor, rutrum eu semper vitae, dapibus sit amet dolor.  

\textsc{Palabras clave:} Nullam, Et Sem Nec, Tellus, Rutrum, Sagittis Ut
\end{abstract}

\afterpage{\blankpage}
\tableofcontents

\newpage
\listoffigures

\afterpage{\blankpage}
\newpage
\listoftables

\afterpage{\blankpage}
\newpage
\listofmyequations

\afterpage{\blankpage}
\lstlistoflistings

\afterpage{\blankpage}
\newpage
\begin{center}
{\huge Índice de Acrónimos}\\[2cm]
\end{center}
\bigskip
\begin{tabular}{ l c l }

\textbf{ASN} & & Aenean Semper Nunc\\
\textbf{QDC} & & Quisque Dapibus Congu\\ 
\textbf{CATS} & & Class Aptent Taciti Sociosqu \\
\textbf{DLM} & & Duis Lacinia Massa \\
\textbf{AIA} & & Aenean Interdum Aliquet \\




\end{tabular}
\afterpage{\blankpage}

\newpage
\begin{center}
{\huge \textit{Agradecimientos}}\\[1.5cm]
\end{center}

Sed id tincidunt magna. Integer metus dui, dictum vitae feugiat eu, fermentum ultrices magna. In sed dignissim massa, eget dictum ex. Aenean metus arcu, condimentum at euismod tempor, consequat quis mi. Orci varius natoque penatibus et magnis dis parturient montes, nascetur ridiculus mus. Suspendisse aliquam nulla tristique nisi euismod, bibendum suscipit risus posuere. Integer porta ipsum tellus, at pharetra felis fringilla vel. \\

Phasellus et dapibus enim, volutpat iaculis augue. Aliquam erat volutpat. Aliquam eu imperdiet enim, sed tincidunt velit. Donec eget porttitor leo. Ut est magna, ullamcorper mollis dignissim vitae, laoreet sed diam. Curabitur eu tempus odio, ac sollicitudin lacus. Suspendisse aliquam sapien a sem aliquam, sit amet ultricies quam condimentum. Sed mollis at urna non commodo. .\\

Praesent eu dapibus dui. Aenean elementum nisi lorem, vitae ultricies felis egestas sed. Cras rhoncus nunc erat, non hendrerit ligula tincidunt id. Integer ut scelerisque nulla, eu commodo enim. 

\afterpage{\blankpage}


\mainmatter 
\pagestyle{thesis}

\chapter{Introducción}
%Escribir un texto de un o dos párrafo(s) máximo de 10 líneas con una introducción al capítulo

Sed interdum tortor at purus lobortis dignissim. Suspendisse sem elit, imperdiet a pharetra at, ultricies eget risus. Nullam eleifend posuere tellus ut fringilla. Maecenas quis leo libero. Pellentesque aliquam sapien eu nunc auctor, a consequat lectus molestie. Proin ullamcorper quam magna, ornare tempus justo rutrum ac. Orci varius natoque penatibus et magnis dis parturient montes, nascetur ridiculus mus. Nullam bibendum consequat massa, sed ornare enim faucibus eget. Sed posuere ac sapien sodales congue.  

\section{Motivación}

Lorem ipsum dolor sit amet, consectetur adipiscing elit. Mauris a tincidunt ligula. Quisque quis sem ut ligula venenatis ultricies auctor eu nisi. In lacinia, diam fermentum cursus scelerisque, velit orci tristique velit, in porttitor nisi purus sed libero. Interdum et malesuada fames ac ante ipsum primis in faucibus. Aliquam sagittis mi vel risus maximus, id volutpat justo tempor. Sed malesuada ante in odio euismod semper sed ut libero. Nulla congue a magna quis vestibulum. Praesent cursus arcu a lectus hendrerit maximus. Nulla lobortis, erat facilisis viverra euismod, lectus nibh rutrum turpis, a dapibus orci nisl ut odio. Duis ultricies auctor orci, tristique lobortis felis sollicitudin sed. 

%Que es lo que te ha motivado para realizar la tesis en esta temática y los aspectos más esenciales que pensabas obtener de ella...

%El habla es uno de los primeros medios que emplea el ser humano para comunicarse. Por ello, desde hace algunos años se investiga el procesamiento del habla con el fin de poder realizar diversas tareas a partir de ella como la transcripción y traducción automática, el reconocimiento de voz, las cuales tienen una gran utilidad. En particular, el reconocimiento del habla tiene diversas aplicaciones como el dictado automático, control por comandos y apoyo para discapacitados.

%Dado que en un inicio no se contaba con la capacidad computacional que se tiene al día de hoy, se emplearon métodos complejos para construir un sistema que realice esta tarea. Estos sistemas contaban con tres modelos: acústico, de lenguaje y de pronunciación, los cuales eran difícil de optimizar. 

%Con el incremento de la capacidad de cómputo y por ende, la introducción del Deep Learning, se logran desarrollar sistemas End-to-End que asignan directamente el audio de entrada a la secuencia de palabras que conforman su transcripción. 


%Este punto podrá ser de 4 páginas máximo. Es una de las partes más importantes que introduce al lector en el trabajo en tu tesis, por lo que debe estar muy bien redactado y estructurado.

\section{Objetivo General}

Vivamus id felis sed sapien egestas finibus. Quisque vulputate sit amet ex rhoncus tempus. Curabitur suscipit eget est eu placerat. Sed dignissim purus vitae eleifend pellentesque. In in posuere erat, eu euismod massa. Ut in nisi elit. Morbi in arcu ut enim semper venenatis lobortis ac odio. Donec varius dui sit amet lacinia placerat. Proin eu justo vitae elit fringilla fringilla eu ut metus. 

Específicamente, los objetivos de este trabajo son:

\begin{itemize}
\item[•] Quisque sed tortor ullamcorper, fermentum mauris vitae, porttitor lacus.
\item[•] Proin aliquam elit ac lectus auctor, nec pharetra massa vulputate.
\item[•] Morbi luctus purus vel dui ornare, vel sollicitudin justo bibendum.

\end{itemize}

Orci varius natoque penatibus et magnis dis parturient montes, nascetur ridiculus mus. Donec mattis porttitor nunc, eget blandit neque sagittis vitae:
\begin{itemize}

\item[•] Sed tempus eros eget augue malesuada, faucibus venenatis quam mattis.
\item[•] Suspendisse eget elit aliquam, semper dolor ac, fermentum purus.
\end{itemize}

\section{Estructura del Seminario}
%Para brindar al lector una idea global del contenido de este trabajo, a continuación se detalla una breve descripción de cada capítulo presente en este seminario de tesis.
\begin{itemize}
\item \textbf{Introducción:} \\
Nullam semper est non tellus vestibulum elementum. In rutrum accumsan orci, non porttitor diam ornare id. Morbi venenatis nibh eleifend velit finibus ornare. Sed sapien libero, finibus ac sem sed, pretium elementum velit. Etiam ullamcorper egestas sapien, in auctor augue accumsan eu. Aliquam sollicitudin semper sem, non rutrum metus. Aenean feugiat ante eu sagittis blandit. Mauris fermentum tempus urna a tincidunt. Aliquam lacinia, risus non consectetur rhoncus, nibh massa rhoncus felis, sit amet tincidunt felis felis ac nulla. Ut mattis semper urna ut hendrerit. Vivamus in ex mattis, finibus metus non, vulputate erat. Integer orci ipsum, aliquet id lorem vitae, rhoncus malesuada ligula. Pellentesque tristique odio in mattis accumsan. 
%En este capítulo introductorio se expone la motivación para desarrollar este trabajo de investigación. Asimismo, se presentan los objetivos y una visión general de la organización del documento.

%En este capítulo introductorio se comenta sobre las motivaciones ...

\item \textbf{Estado del Arte:} \\
nterdum et malesuada fames ac ante ipsum primis in faucibus. Nunc sollicitudin, lorem rhoncus sollicitudin fringilla, nibh turpis consectetur nisl, sed iaculis justo massa at risus. Pellentesque volutpat euismod mi, nec gravida justo dictum id. Etiam iaculis tortor dolor, eget egestas eros luctus non. Quisque consectetur egestas arcu vitae blandit. Etiam pretium tortor ullamcorper nibh porta, sit amet egestas justo auctor. Aliquam gravida venenatis metus id iaculis. 
%La base teórica sobre la cual se sustenta el presente trabajo se detalla en este apartado. Se exponen, además, conceptos con respecto a la navegación de un robot móvil, la utilización de sensores y cámaras para identificar el ambiente donde se movilizará el robot así como los métodos para su localización y algoritmos de enrutamiento.
%Este apartado presenta la forma en como se genera el habla así como también se muestran conceptos básicos de Deep Learning. Finalmente, se expone el problema que busca resolver este trabajo.

\item \textbf{Herramientas:} \\
Pellentesque consequat auctor dolor quis viverra. Aenean a felis pulvinar, imperdiet urna ac, fermentum lorem. Integer porta turpis nec ex fringilla, a posuere tortor tristique. Nullam pulvinar augue nec facilisis ultricies. Sed vel fermentum neque. Fusce erat metus, varius et gravida nec, aliquam a ligula. Pellentesque habitant morbi tristique senectus et netus et malesuada fames ac turpis egestas. 
%El capítulo se centra en la descripción de librerías empleadas para el procesamiento de los datos y la creación de las arquitecturas de red.

%En este capítulo se describe el sistema empleado, sus componentes electrónicos y el framework utilizado para la comunicación entre la tarjeta controladora y la microcomputadora. 

\item \textbf{Análisis de Datos de Entorno} \\
%En este capítulo se detalla la base de datos empleada, las transformaciones que se aplicaron a los datos y el aumento de los datos.
Aliquam ex felis, malesuada vel blandit ac, finibus vitae odio. Nulla facilisi. Etiam euismod vehicula risus, id aliquet nibh. Quisque auctor venenatis molestie.  

\item \textbf{Desarrollo de la investigación:} \\
%En este capítulo se detalla el problema que busca resolver el presente seminario y la propuesta de solución.
%En este capítulo se describe la ejecución del 
Nulla vitae elit nec ligula lobortis aliquam. Aliquam lobortis iaculis nibh semper bibendum. Aliquam feugiat mi et velit laoreet, ac blandit turpis blandit. Phasellus posuere, nisl a vestibulum finibus, nisi urna luctus magna, quis elementum nunc risus sit amet sem. Fusce non ligula interdum, mollis quam at, iaculis lorem. Maecenas consequat volutpat nulla ac scelerisque. 

\item \textbf{Análisis de Resultados:} \\
Interdum et malesuada fames ac ante ipsum primis in faucibus. Ut cursus tempor consectetur. Praesent eleifend arcu ac mattis mollis. Morbi at lobortis nunc.
%En este capítulo se muestran los resultados experimentales obtenidos y se comparan con los resultados esperados.
%En este apartado se analiza la exactitud del sistema ASR desarrollado y se realiza una comparación entre los diversos parámetros utilizados.

\item \textbf{Conclusiones y Trabajo Futuro:}\\
Curabitur a blandit lectus. Etiam dui mi, interdum sed gravida ac, faucibus vel justo. Sed tellus erat, rhoncus at nibh sed, convallis finibus leo. Sed commodo nec arcu nec sodales. 
%El capítulo se centra en exponer las conclusiones obtenidas de este trabajo. Adicionalmente, se proponen trabajos a futuro que puedan ser implementados con el fin de mejorar la navegación.
%El capítulo se centra en exponer las conclusiones de acuerdo a los objetivos propuestos para este Seminario de Tesis II. Adicionalmente, se proponen trabajos a futuro que puedan ser implementados con el fin de mejorar el sistema ASR.
%para la implementación de el servidor y la aplicación web del sistema.
\end{itemize}

%NOTA: RECUERDE QUE ES COMO UN LIBRO TODO CAPÍTULO NUEVO COMENZARÁ CON PÁGINA IMPAR NUNCA PAR

\afterpage{\blankpage}



%\newpage
%$\ $
%\thispagestyle{empty} % para que no se numere esta pagina


\chapter{Estado del Arte}
Suspendisse potenti. Maecenas eu diam sit amet metus vulputate commodo. Praesent dapibus vestibulum nulla. Etiam nec pretium ante. Praesent velit metus, lobortis eu ipsum sed, elementum interdum sem. Donec sed porttitor libero, a pretium lectus. Etiam commodo, ipsum eu venenatis commodo, eros nisl posuere purus, a ullamcorper magna nisl et justo. Donec ornare nisi sed sapien tempus molestie. In eget mi nisl. Nunc eget sapien lectus. Nulla egestas vehicula ante ut eleifend. Aenean lobortis, massa in porta rhoncus, enim sapien congue ex, quis venenatis ante elit id lacus. Nunc tellus nibh, ultrices non mauris non, varius laoreet lorem. 

\section{El Entorno y el Agente}
Proin aliquam vestibulum nibh eget aliquet. Nunc mattis, nisi sed maximus mattis, sem elit pretium nibh, et finibus erat libero id justo. Integer luctus, ex in congue elementum, turpis quam commodo sapien, at porta enim lectus vel purus. Proin a ultricies elit. Nullam nibh quam, vehicula et vehicula eu, consequat eu eros. 
\subsection{Entorno: Air Hockey}
Maecenas arcu ligula, semper ac ipsum ut, tempor ultrices quam. Aenean id ultricies mauris. Ut mattis enim a arcu sollicitudin, pharetra vehicula odio iaculis.  \cite{usaa 2022}.Maecenas vitae ultricies dui, vitae dignissim dolor. Sed pulvinar ullamcorper odio nec sagittis. Nullam congue velit sem, vitae fermentum felis vestibulum ut.  figura \ref{fig:air_hockey}. Cras ipsum justo, venenatis sit amet est laoreet, tempus facilisis leo. Integer blandit leo massa, a ullamcorper nisi semper quis. Fusce dapibus sit amet justo a luctus. Fusce nunc diam, commodo a porta et, elementum nec purus. Interdum et malesuada fames ac ante ipsum primis in faucibus. \cite{airhockey 1975},Ut pellentesque, augue pretium congue pretium, lorem neque vehicula mauris, interdum fermentum arcu leo viverra ipsum. Aenean pulvinar sem ligula, id lacinia erat volutpat in.  
%Principalmente en nuestra investigación se plantea poder reemplazar los jugadores humanos por robots Scara, todo ello de forma simulada en el ambiente controlado PyBullet.
\begin{figure}[H]
\centering
\includegraphics[width=0.8\textwidth]{Figures/airhockey_demostracion.jpg}
\caption[Air Hockey]
	{Vivamus vitae imperdiet tellus. In finibus tincidunt tempus. Phasellus porta mi eget ante imperdiet, in tincidunt sem pharetra. Vestibulum tincidunt nisi non porttitor auctor \cite{airhockey 1975}.}
\label{fig:air_hockey}
\end{figure}

\subsection{Agente: Robot Scara}
In vel massa gravida, semper sapien ac, vulputate ipsum. Pellentesque a libero nec libero sodales egestas. Integer in massa sodales, cursus nibh sit amet, tristique tortor. Nunc non gravida est. Nam imperdiet ultricies purus sed eleifend. Duis a rhoncus tortor, eget malesuada ex. Pellentesque habitant morbi tristique senectus et netus et malesuada fames ac turpis egestas. Nunc at tellus ac magna vestibulum consectetur luctus nec tellus. Cras eu scelerisque mauris, non suscipit velit. Nam malesuada dui eu odio laoreet cursus. Curabitur vulputate lacus et lacus interdum, a pretium purus rhoncus \cite{Reyes 2022}.

\begin{figure}[H]
\centering
\includegraphics[width=0.45\textwidth]{Figures/scara.png}
\caption[Robot Scara]
	{Phasellus gravida dui dolor \cite{airhockey 1975}.}
\label{fig:filtro}
\end{figure}

\subsubsection{Modelo de Cinemática Directa}
Sed sollicitudin velit lorem, ac mollis felis tincidunt vel. Pellentesque tincidunt diam neque, sit amet aliquam sapien convallis eget. Duis in pellentesque turpis. Vestibulum mi erat, tincidunt at velit eu, tempor sodales augue. Vestibulum purus turpis, congue non lacus in, condimentum convallis elit. Phasellus sit amet nisi aliquam, imperdiet orci sed, lacinia nibh. Nullam nec tempus ipsum. Nullam sed nisi mollis, vulputate nunc sit amet, condimentum leo figura 2.3.
\begin{figure}[H]
\centering
\includegraphics[width=0.8\textwidth]{Figures/cinematica_geometrica.png}
\caption[Cinemática Directa Robot Scara]
	{ In mauris nisl, porta eleifend viverra vitae \cite{Reyes 2022}.}
\label{fig:filtro}
\end{figure}

\section{Marco teórico}

\subsection{Aprendizaje por Refuerzo}
In hac habitasse platea dictumst. In eget vehicula massa, quis ornare orci. Interdum et malesuada fames ac ante ipsum primis in faucibus. Duis sit amet quam nisi \cite{Sutton 2022}:
\begin{itemize}
    \item \textbf{Agente.}
    Nulla nec augue vitae nulla elementum fringilla. Pellentesque at ultrices ipsum. Ut ultricies volutpat tellus, a rutrum ligula bibendum pulvinar.
    \item \textbf{Entorno}
    Curabitur sagittis eu neque id rutrum. Ut non leo ac ante lobortis volutpat nec non metus. Fusce varius cursus tellus at ultrices. Donec at pharetra sapien.
    \item \textbf{Estado}
    Nam tempor ullamcorper tincidunt. Phasellus luctus eu massa nec varius. Praesent rhoncus, libero at euismod mollis, urna lorem suscipit est, vel finibus nisl magna sit amet nulla.
    \item \textbf{Acción}
    In a nibh tincidunt, commodo lectus et, placerat nisl. Suspendisse ut justo eu odio blandit maximus. Sed in dignissim justo. Aliquam hendrerit purus vitae fringilla varius.
    \item \textbf{Recompensa}
    Cras at nisi et orci sollicitudin commodo. Aenean lacinia ipsum vel fringilla egestas. Mauris at hendrerit magna, in consequat justo. 
\end{itemize}
Praesent velit erat, varius eu imperdiet ac, euismod sed nunc. Aliquam dignissim sem ante, et pulvinar tortor feugiat lobortis. \cite{L. Lozano} ilustrado en la figura 2.4.
    
\begin{figure}[H]
\centering
\includegraphics[width=0.40\textwidth]{Figures/esquema_RL.png}
\caption[Aprendizaje por Refuerzo, ciclo de operación]
	{In hac habitasse platea dictumst. Pellentesque sit amet \cite{Reyes 2022}.}
\label{fig:filtro}
\end{figure}





\subsection{Redes Neuronales}

\begin{itemize}
    \item \textbf{Cras}
    Vivamus faucibus arcu nisi, vel vulputate risus suscipit eget. Proin a consequat lacus. Ut bibendum laoreet aliquet. Interdum et malesuada fames ac ante ipsum primis in faucibus. figura 2.5.
    \begin{itemize}
        \item
    \end{itemize}
    \item \textbf{Capa}
    
    \item \textbf{Red}
    
    \item \textbf{Sistema Neuronal}
    
    
\end{itemize}

\begin{figure}[H]
\centering
\includegraphics[width=0.70\textwidth]{Figures/Red_neuronal.png}
\caption[Red Neuronal, componentes]
	{Nullam arcu urna, porttitor id sagittis quis, hendrerit in magna \cite{Reyes 2022}.}
\label{fig:filtro}
\end{figure}


\subsection{Mauris commodo lobortiso}

\subsection{Pellentesque nec nisi id metus varius tristi}

\subsection{Aliquam sit amet justo}

\subsection{Sed efficitur tortor at orci}

\subsection{Modelo de Actor-Crico}

\subsection{Curabitur posuere diam velit, et ultricies nunc eleifend nec}

\subsection{Donec eu dui felis}

\subsection{Pellentesque condimentum lectus}


\section{Artículos Relacionados}





















\begin{comment}
\section{Problema A}

\section{Concepto A}


\section{Generación del habla}

\subsection{Modelo Source-Filter}
Un filtro acústico atenúa selectivamente (reduce en intensidad) ciertas frecuencias y permite que otras frecuencias pasen relativamente sin atenuar. La teoría del source-filter describe la producción del habla como un proceso de dos etapas \cite{book:soundLanguage}:
\begin{itemize}
    \item La generación de una fuente de sonido, con su propia forma espectral y estructura espectral fina.
    \item El filtrado del espectro por las propiedades resonantes del tracto vocal.
\end{itemize}

La mayor parte del filtrado de un espectro de fuente se lleva a cabo por la parte del tracto vocal anterior a la fuente de sonido. Por ejemplo, en el caso de una fuente glotal, el filtro es todo el tracto vocal supraglótico.





El filtro del tracto vocal siempre incluye alguna parte de la cavidad oral y también, opcionalmente, puede incluir la cavidad nasal. Para la producción de un sonido nasal, se
bloquea el pasaje oral y se levanta el velo que normalmente bloquea el pasaje nasal.



Las fuentes de sonido pueden ser periódicas o aperiódicas. Las fuentes de sonido glotal
pueden ser:

\begin{enumerate}
    \item \textbf{Periódicas (sonoras).} Por ejemplo, la producción de la /a/, la glotis se abre y se cierra periódicamente.
    \item \textbf{Aperiódicas (susurros y /h/).} Durante la producción de sonidos sordos, la glotis no vibra y está abierta. Por ejemplo, en la producción de /s/.
    \item \textbf{Mixtas (por ejemplo, voz entrecortada).} Por ejemplo, en la producción de /z/.
\end{enumerate}

\section{Reconocimiento automático del habla (ASR)}
Un sistema ASR produce la secuencia de palabras más probable dada una señal de voz
entrante.

\subsection{Modelos para el reconocimiento de voz de vocabulario extenso (LVCSR)}



\subsubsection{A. Tradicional: Modelo basado en HMM}
Se divide en tres partes modelos, los cuales son independientes entre sí y donde cada uno desempeña un papel diferente:

\begin{itemize}
\item \textbf{Modelo Acústico.} Determina el mapeo entre la entrada del habla y la secuencia de características. Existen dos tipos HMM-GMM y HMM-DNN. El HMM se utiliza principalmente para realizar deformaciones de tiempo dinámicas a nivel de frame. GMM y DNN se utilizan para calcular la probabilidad de emisión de los estados ocultos del HMM.

\begin{itemize}
    \item \textbf{Modelo HMM-GMM.} La probabilidad de observación está representada generalmente por GMM (Gaussian Mixture Model).
    \item \textbf{Modelo HMM-DNN.} La distribución de probabilidad posterior del estado oculto se puede calcular mediante DNN. Proporciona mejores resultados.
    
\end{itemize}

\item \textbf{Modelo de Pronunciación.} Define el mapeo entre fonemas (o subfonemas) y grafemas. Es la conexión entre la secuencia acústica y la secuencia del lenguaje.

\item \textbf{Modelo de Lenguaje.} Es el mapeo la secuencia de caracteres para una transcripción final fluida. Predice la probabilidad de que palabras específicas aparezcan una
tras otra en un idioma determinado. Los reconocedores típicos utilizan modelos de lenguaje n-gram. Un n-grama contiene la probabilidad previa de la aparición de una
palabra (unigrama) o de una secuencia de palabras (bigrama,
trigrama, etc.):

\begin{itemize}
    \item Probabilidad de unigrama: $P(w_i)$
    \item Probabilidad de bigrama: $P(w_i|w_{i-1})$
    \item Probabilidad de n-grama: $P(w_n|w_{n-1},w_{n-2},\dots,w_1)$
\end{itemize}

\end{itemize}

Limitaciones de los modelos HMM:

\begin{itemize}
    \item El proceso de formación es complejo y difícil de optimizar globalmente.
    \item Utiliza supuestos de independencia condicional dentro de HMM. Esto no coincide con la situación real de LVCSR.
\end{itemize}

\subsubsection{B. Empleando Deep Learning: Modelo End-to-End}
Este modelo es un sistema que asigna directamente la secuencia de audio de entrada a la secuencia de palabras u otros grafemas. 


Estructura Funcional:

\begin{itemize}
    \item \textbf{Encoder.} Asigna la secuencia de entrada de voz a la secuencia de características.
    \item \textbf{Aligner.} Encuentra la alineación entre la secuencia de características y el lenguaje.
    \item \textbf{Decoder.} Decodifica el resultado de identificación final.
\end{itemize}

Ventajas:

\begin{itemize}
    \item Varios módulos se fusionan en una red para la formación conjunta, es decir, no es necesario diseñar muchos módulos para realizar el mapeo entre varios estados intermedios.
    \item A diferencia de los modelos basados en HMM, el modelo End-to-End no presenta una representación interna para la pronunciación de una cadena de caracteres.
\end{itemize}

Según sus implementaciones de alineamiento suave, este modelo se divide en tres categorías:

\begin{itemize}
    \item \textbf{Basado en CTC.} Primero enumera todas las posibles alineaciones dura. Luego, logra una alineación suave agregando estas alineaciones dura. CTC asume que las etiquetas de salida son independientes entre sí al enumerar alineaciones duras.
    \item \textbf{Transductor RNN.} Enumera todas las posibles alineaciones duras y luego las agrega para una alineación suave. Pero a diferencia de CTC, el transductor RNN no hace suposiciones independientes sobre las etiquetas al enumerar alineaciones duras. Por lo tanto, es diferente de CTC en términos de definición de ruta y cálculo de probabilidad.
    \item \textbf{Basado en la atención.} No enumera todas las posibles alineaciones duras, sino que utiliza un mecanismo de atención para calcular directamente la información de alineación suave entre los datos de entrada y la etiqueta de salida.
\end{itemize}

\section{Deep Learning}

En esta sección presentamos los conceptos generales empleados en la construcción de un sistema ASR.

\subsection{Redes Neuronales Recurrentes (RNNs)}
Las redes neuronales recurrentes (RNN) se originan de introducir conexiones recurrentes dentro de las capas de una FNN (Feed Forward Neural Network) de tal forma que cada paso de tiempo tiene como una de sus entradas el estado anterior$t$. Esta característica permite que tengan memoria. 

Dada una secuencia de entrada $x(t) = (x_1, \dots, x_T)$ en el paso de tiempo actual $t$, las RNN obtienen el estado oculto, $h_t$, utilizando el estado oculto anterior, $h_t − 1$, y generan una secuencia vectorial $y = (y_1, \dots, y_T)$\cite{11_2001.00378}. El cálculo de estos elementos se muestran en \ref{eq:estado_oculto} y \ref{eq:yt}
\begin{equation}
    h_t = H(W_{xh}x_t + W{hh}h_{t−1} + b_h
\end{equation}
\label{eq:estado_oculto}
\myequations{Estado oculto}

\begin{equation}
    y_t = (W{xh}x_t + b_y)
\end{equation}
\label{eq:yt}
\myequations{Secuencia Vectorial}

donde W es las matriz de ponderación (es decir, $W_{xh}$ es una matriz de ponderación de una capa oculta de entrada), $b$ es el vector de sesgo y $H$ denota la función de capa oculta. 

Lastimosamente este tipo de redes presenta el problema de gradiente de desaparición (vanishing gradient), en el cual la contribución de estados anteriores es cada vez menor con el paso del tiempo. Las aquitecturas de \textbf{memoria a corto / largo plazo (LSTM}) y \textbf{las unidades recurrentes controladas (GRU)} resuelven este problema al manejar mecanimos de activación para agregar y olvidar información determinada.

Existen dos tipos de RNNs:

\begin{itemize}
    \item \textbf{RNN Unidireccional.} Son las que se han presentado hasta ahora en la que se emplean estados anteriores.
    \item \textbf{RNN Bidireccional.} Emplea estados anteriores y posteriores, lo cual mejora la exactitud de la red. Sin embargo, es de alta latencia ya que es necesario procesar la secuencia completa.
\end{itemize}

\subsection{Métricas}
El rendimiento de un sistema ASR se mide comparando las transcripciones hipotéticas y las transcripciones de referencia.

\subsubsection{Tasa de error de palabra (WER)}
Es la métrica más utilizada. Las dos secuencias de palabras se alinean
primero usando un algoritmo de alineación de cadenas basado en
programación dinámica. Después de la alineación, se determina el número
de eliminaciones (D), sustituciones (S) e inserciones (I). Las eliminaciones,
sustituciones e inserciones se consideran errores, y el WER se calcula por la
tasa del número de errores al número de palabras (N) en la referencia.

\begin{equation}
    WER = \frac{I+D+S}{N}*100\%
\end{equation}
\label{eq:wer}
\myequations{WER}

\subsubsection{Tasa de error de caracter (CER)}
Calcula el porcentaje de oraciones con al menos un error.
\end{comment}
%\subsection{Trabajo Relacionado}
%En la literatura se puede encontrar diversas referencias a sistemas y/o métodos de ASR. A continuación, se presentan algunos de ellos. \\

%La arquitectura presentada en \cite{01_1512.02595} consta de capas convulaciones invariantes de una o dos dimensiones (time-only domain o time-and-frequency domain) para reducir la cantidad de pasos de tiempo. Seguida de estas, se encuentran un conjunto de capas recurrentes LSTM unidireccionales y finalmente, seguidas de capas completamente conectadas. Para el entrenamiento de la red de emplea la función de pérdida CTC y el aprendizaje curricular (curriculum learning) SortaGrad que toma como heurística de dificultad a la longitud de las declaraciones, es decir, entra las declaraciones de menor a mayor longitud. Adicionalmente, emplea un modelo de lenguaje de 5-grama para el idioma inglés.

%En modelo presentado en \cite{04_2006.11477} está compuesto de un encoder de características convolucional multicapa que toma como entrada el audio y tiene como salida representaciones del habla que posteriormente son discretizadas para representar los targets.



\afterpage{\blankpage}
%\newpage
$\ $
%\thispagestyle{empty} % para que no se numere esta pagina


\chapter{Herramientas}
En este capítulo se describe las herramientas usadas en esta investigación y la dividiremos en dos secciones, la primera describirá herramientas para la construcción del entorno personalizado AIR HOCKEY y el agente robot Scara simulado, finalmente en la segunda se mostrarán herramientas relacionadas para poder construir la parte lógica que realice el control del agente.

\section{Herramientas de Construcción de entorno Personalizado AIR HOCKEY}
     \subsection{URDF - United Robotics Description Format}
     Es un archivo escrito en formato en el marco de gramática XML y nos permite mediante etiquetas definir la forma visual de los componentes u objetos que necesitamos en nuestro entorno o robot, este formato nos permite poder crear enlaces entre componentes mediante la arquitectura de arbol, donde las ramas son las uniones y los nodos son los componentes. Además de esto las propiedades tales como color, friccion, inercia, colision de un objeto especifico deben ser incluidas en este archivo para poder ser cargados posteriormente por un simulador.
    \subsection{PyBullet}
    Es el simulador escogido en esta investigación y contiene un Motor de Físicas que pretende simular un entorno Físico Real, la razón de usar esta herramienta es debido a que podemos usar diversas funciones que nos permiten personalizar condiciones físicas, tales como inercia, fricción de superficies, gravedad, colisiones, entre otras. Por poner un ejemplo claro, en el juego real de AIR HOCKEY el disco de juego se mueve libremente gracias a una delgada película de aire que lo mantiene flotando y sin fricción, para abordar esta peculiaridad, podemos usar la propiedad Fricción Lateral del URDF cargado del tablero y setearlo a 0, de esta forma simplificamos la necesidad de modelar una película de aire y conseguir el mismo fenómeno.
   
\section{Herramientas de Construcción de Lógica de Agente}
    \subsection{Pytorch}
    Nos permite construir nuestra Arquitectura de Red Neuronal, será util para poder crear las redes neuronales de Críticos Gemelos.


%\afterpage{\blankpage}
%\newpage
$\ $
%\thispagestyle{empty} % para que no se numere esta pagina
\chapter{Análisis}

\section{Ut vitae lacinia}


\section{Aenean lacinia cursus}

% intro al cap
\section{Donec quis erat}
Vivamus sit amet nulla quis lectus facilisis ultricies ac a est. Donec fringilla pellentesque mi. Mauris at augue eu sapien maximus commodo in eget eros. Suspendisse potenti. Donec nec diam interdum, aliquet leo sit amet, imperdiet ex. Cras a laoreet lacus, eu tempor est. Curabitur vitae mi placerat, efficitur nulla tempus, mattis nisl. Nulla facilisi. 
\begin{table}[H]
\centering
\begin{tabular}{| c | c | c|} 
 \hline
 \textbf{Acción} & \textbf{Recompensa} & \textbf{Inconvenientes}\\ [0.3ex] 
 \hline 
 Mazo Golpea disco &  &   \\ 
 \hline
 Mazo delante de Disco & &  \\
 \hline
 Disco Anota en la portería contraria &  &  \\
 \hline 
 Disco entra en la portería del agente &  &   \\ %[1ex] 
 \hline 
 Disco se mueve hacia porteria contraria &  &   \\ %[1ex] 
 \hline
\end{tabular}
\caption[Movimientos Básicos Programados]
	{Proin blandit risus laoreet nisl pellentesque, eget aliquam nibh facilisis}
\label{table:movimientos_basicos}
\end{table}


\afterpage{\blankpage}
%\newpage
$\ $
%\thispagestyle{empty} % para que no se numere esta pagina
%\thispagestyle{empty} % para que no se numere esta pagina
\chapter{Desarrollo de la investigación}



\section{Concepto A en el Desarrollo de la Investigación con el input}

\lipsum[1-1]

\dots se detallan en el Apéndice~\ref{AppendA}.

\section{Teorías C  en el Desarrollo de la Investigación con el input}

\lipsum[1-1]

\dots en la Tabla \ref{table:movimientos_basicos}.

\begin{table}[H]
\centering
\begin{tabular}{| c | c | c|} 
 \hline
 \textbf{Modo} & \textbf{Articulaciones} & \textbf{Ruedas}\\ [0.3ex] 
 \hline \hline
 Avanza & $J_1=J_4=512$, $J_2=J_3=-512$ & $W_1=W_3=-v$, $W_2=W_4=v$  \\ 
 \hline
 Retrocede & $J_1=J_4=512$, $J_2=J_3=-512$ & $W_1=W_3=v$, $W_2=W_4=-v$  \\
 \hline \hline 
 Giro CCW & $J_1=J_2=J_3=J_4=0$ & $W_1=W_2=W_3=W_4=v$ \\
 \hline 
 Giro CW & $J_1=J_2=J_3=J_4=0$ & $W_1=W_2=W_3=W_4=-v$  \\ %[1ex] 
 \hline \hline 
 Derecha & $J_1=J_4=-512$, $J_2=J_3=512$ & $W_1=W_2=-v$, $W_3=W_4=v$  \\ %[1ex] 
 \hline
 Izquierda & $J_1=J_4=-512$, $J_2=J_3=512$ & $W_1=W_2=v$, $W_3=W_4=-v$  \\
 \hline
\end{tabular}
\caption[Movimientos Básicos Programados]
	{Movimientos Básicos Programados. $J_i$ es el valor de la posición para el servomotor que controla la articulación de la i-ésima extremidad, y donde $W_i$ es el valor de la velocidad de la rueda de la i-ésima extremidad.  Fuente: Elaboración propia.}
\label{table:movimientos_basicos}
\end{table}

\section{Metodología E  en el Desarrollo de la Investigación con el input}
\lipsum[1-1]

\begin{lstlisting}[language=Python, caption=Ejemplo en Python, label={lst:alg1}, numbers=left]
import numpy as np
    
def sum(x,y):
    return x + y
\end{lstlisting}

\lipsum[1-1]

\dots el Algoritmo \ref{lst:alg1}.

%\afterpage{\blankpage}
%\newpage
$\ $
%\thispagestyle{empty} % para que no se numere esta pagina
\chapter{Análisis de Resultados}

% intro al cap
\section{Comparación}
\section{Mejora}
\section{Descripción}
\section{Estadísticas}
\section{Validación}

\afterpage{\blankpage}
%\newpage
$\ $
%\thispagestyle{empty} % para que no se numere esta pagina
\chapter{Conclusiones y Trabajo Futuro}
%intro al cap

%ESTE CAPÍTULO ES UNO DE LOS MÁS IMPORTANTES, POR NO DECIR EL QUE MÁS. EN ÉL, EL JURADO VA A TENER CLARO QUÉ HA APRENDIDO EL ALUMNO Y CÓMO LO HA DESARROLLADO, LOS PROBLEMAS QUE HAN SURGIDO Y COMO LOS HA SOLUCIONADO...
%ADEMÁS DE QUE EL ALUMNO DEJARÁ CLARO QUE SE HA ESPECIALIZADO EN LA TEMÁTICA Y DEJARÁ EN ESCRITO TODO LO APRENDIDO Y COMO CONTINUARÁ CON LA TEMÁTICA EN POSTERIORES ESTUDIOS DEL MISMO TEMA

\section{Conclusiones}


\section{Trabajo Futuro}

\subsection*{Trabajo Futuro 1}

\subsection*{Trabajo Futuro 2}

%\afterpage{\blankpage}

\begin{thebibliography}{0}
 
 \bibitem{usaa 2022} Official United States Air Hockey Association. Recuperado el 11 de Octubre de 2022 https://airhockeypros.com/fundamentals.html
 
 \bibitem{airhockey 1975} G. M. Kaler, "Build this tabletop hockey game for family fun", Popular Mechanics, vol. 144, n.º 5, p. 188, 1975. 
 
  \bibitem{Reyes 2022} F. Reyes, Robótica - control de robots manipuladores, 1ª ed, México D.F: AlfaOmega, 2011, pp. 236-238
  
  \bibitem{Kelly 2003} R. Kelly, V. Santibañez, Control de movimiento de robots manipuladores. Madrid: Pearson Educación, 2003, pp. 285 - 297.
  
  \bibitem{Sutton 2022} R. Sutton, Reinforcement learning: an introduction, Cambridge, Massachusetts: MIT Press, 1998.
  
  \bibitem{L. Lozano} L. Lozano, Trabajo de fin de grado en Matemáticas: Aprendizaje por Refuerzo,Elementos básicos y algoritmos. España, Universidad de Zaragoza.

\end{thebibliography}
\afterpage{\blankpage}

\appendix
% Appendix A

\chapter{Este es un apéndice} % Main appendix title

\label{AppendA} % For referencing this appendix elsewhere, use \ref{AppendixA}

% Appendix B

\chapter{Este es otro apéndice} % Main appendix title

\label{AppendB} % For referencing this appendix elsewhere, use \ref{AppendixB}

\section{Sección del apéndice}

% Appendix C

\chapter{Este es mi último apéndice} % Main appendix title

\label{AppendC} % For referencing this appendix elsewhere, use \ref{AppendixC}

\section{Primera Sección}

\section{Segunda Sección}

\end{document}